\subsection{The algebraic-category obstruction in the four-pointed case}

\begin{proposition}[Refuted: finite-type algebraicity of the natural cut component]
\label{A1-cut-not-algebraic}
Let $\mathcal D$ be a connected component of the moduli space of
four-labelled punctured spheres equipped with one essential
unoriented cut isotopy class, represented by its hyperbolic
geodesic. Give $\mathcal D$ its natural complex structure: forgetting
the locally constant cut class is a local biholomorphism to
$M_{0,4}^{\mathrm{an}}$.
Then
\[
 \mathcal D\simeq
 \mathbb H/\langle\tau\mapsto\tau+2\rangle
 \simeq\{q\in\C:0<|q|<1\}.
\]
In particular $\mathcal D$ is not the analytification of a smooth
connected finite-type algebraic curve over $\C$. Thus the natural
geodesic-cut source cannot be treated as a finite-type algebraic
moduli curve, and it cannot be algebraically isomorphic to the
corresponding unrestricted nodal-tensor stratum $\C^\times$.
This is a limitation of that specified algebraic interpretation;
it does not infer nonholomorphicity of the surgery map from its
real-analytic description.
\end{proposition}
\begin{proof}
We give the passage from the elliptic double cover explicitly.
The double cover of a four-pointed sphere branched at its four
markings has genus one by Riemann--Hurwitz. Choose one ramification
point as origin. The other three are its nonzero two-torsion
points. Conversely an elliptic curve divided by $z\mapsto-z$
recovers the sphere and these four markings. Markings lift in two
ways, differing by $-1$. Thus the effective marked parameter is
the period ratio $\tau\in\mathbb H$ for
$E_\tau=\C/(\Z+\tau\Z)$. This is a holomorphic identification
of the Teichm\"uller spaces: locally the double-cover family has
equation $y^2=x(x-1)(x-\lambda)$, and the ratio of the two marked
periods of $dx/y$ varies holomorphically; the inverse is the
quotient elliptic family. The torus parameter itself is
\cite[Example 2.3, p.~12]{AHubbardKoch2014}.

The lifting and curve correspondence are proved in
\cite[Propositions 2.6--2.7, pp.~55--57]{FarbMargalit2012A567}.
Here is the additional restriction needed for labelled markings.
Fixing the chosen origin removes the two-torsion translations;
fixing all four labels requires the homology matrix to be the
identity modulo two. The effective pure group is consequently
$\Gamma(2)/\{\pm I\}$. Under the double cover an essential sphere
curve lifts to two parallel primitive torus curves. Choose a
homology basis in which their common unoriented class is
$\pm(1,0)$. Its stabilizer in the pure group consists exactly of
the matrices
\[
 \pm\begin{pmatrix}1&2m\\0&1\end{pmatrix},\qquad m\in\Z.
\]
Indeed preserving that primitive line forces $c=0$ and $a=d=\pm1$,
and congruence modulo two forces $b$ even. A Dehn twist downstairs
lifts to the twists around the two parallel curves, and is the
generator with $|m|=1$. Therefore forgetting the full marking but
retaining the cut gives precisely the claimed cyclic quotient.
The exponential $q=\exp(\pi i\tau)$ identifies it biholomorphically
with the punctured unit disc. This also agrees with the cyclic
parabolic quotient in \cite[Example 7.2, pp.~25--26]{AHubbardKoch2014},
with period two instead of period one. Since an automorphism of a
sphere fixing its four labelled points is the identity, no residual
orbifold stabilizer changes this calculation.

Suppose now that this punctured disc were $X^{\mathrm{an}}$ for a
smooth connected finite-type algebraic curve $X/\C$. Its smooth
projective completion has some genus $g$, and its complement is a
finite nonempty set of $r$ points; nonemptiness follows because
the punctured disc is not compact. The fundamental group of the
resulting punctured surface is free of rank $2g+r-1$. Since the
punctured disc has fundamental group $\Z$, we have $2g+r=2$,
and $r\ge1$ forces $g=0,r=2$. Thus $X\simeq\mathbf P^1\setminus
\{0,\infty\}=\C^\times$. But $\C^\times$ has universal cover
$\C$, whereas the punctured disc has universal cover $\mathbb H$;
these simply connected Riemann surfaces are not biholomorphic.
For example, after identifying $\mathbb H$ with $\D$, a
biholomorphism $\C\to\mathbb H$ would contradict Liouville's
theorem. This proves the contradiction.

The fixed labelled one-node normalization in the target has no
moduli or automorphisms, and its nonzero tangent tensor is an
arbitrary element of a one-dimensional complex vector space.
Hence its unrestricted target stratum is exactly $\C^\times$.
The stated algebraic isomorphism is therefore impossible for the
natural source just analyzed, independently of the cap-size
obstruction and without any assertion about whether a particular
real-analytic formula is holomorphic.
\end{proof}

Both the punctured disc and $\C^\times$ have fundamental group
$\Z$. The local real tangent isomorphism and the fundamental-group
isomorphism in Theorem~\ref{A1-graph-range} are thus fully
compatible with the failure of a global isomorphism onto the
unrestricted target. They cannot substitute for the range
condition or for the algebraic-category hypothesis.
